\documentclass[aspectratio=169]{beamer}

% Theme and color scheme
\usetheme{default}
\usecolortheme{default}

% Remove navigation symbols
\setbeamertemplate{navigation symbols}{}

% Simple color scheme for printing
\definecolor{darkblue}{RGB}{25,25,112}
\setbeamercolor{title}{fg=darkblue}
\setbeamercolor{frametitle}{fg=darkblue}
\setbeamercolor{structure}{fg=darkblue}

% Simple footline
\setbeamertemplate{footline}[frame number]

% Packages
\usepackage{amsmath}
\usepackage{amsfonts}
\usepackage{amssymb}
\usepackage{graphicx}
\usepackage{tikz}
\usepackage{booktabs}
\usepackage{array}
\usepackage{listings}
\usepackage{xcolor}

% Code listing setup
\lstdefinestyle{chisel}{
    language=Scala,
    basicstyle=\ttfamily\footnotesize,
    keywordstyle=\color{blue}\bfseries,
    commentstyle=\color{green!60!black},
    stringstyle=\color{red},
    numbers=left,
    numberstyle=\tiny\color{gray},
    stepnumber=1,
    numbersep=5pt,
    backgroundcolor=\color{gray!10},
    frame=single,
    rulecolor=\color{black!30},
    breaklines=true,
    breakatwhitespace=false,
    tabsize=2,
    captionpos=b
}

\lstdefinestyle{verilog}{
    language=Verilog,
    basicstyle=\ttfamily\footnotesize,
    keywordstyle=\color{blue}\bfseries,
    commentstyle=\color{green!60!black},
    stringstyle=\color{red},
    numbers=left,
    numberstyle=\tiny\color{gray},
    stepnumber=1,
    numbersep=5pt,
    backgroundcolor=\color{gray!10},
    frame=single,
    rulecolor=\color{black!30},
    breaklines=true,
    breakatwhitespace=false,
    tabsize=2,
    captionpos=b
}

% Set graphics path for images (relative to slides directory)
\graphicspath{{../Images/}}

% Title information
\title[Appendix A01: Hardware Description Languages]{Structured Computer Architecture\\Appendix A01: Hardware Description Languages\\Chisel, SystemVerilog, and Comparisons}
\author{Based on Structured Computer Architecture: An Order-Based Approach}
\date{}

% Custom command for highlighting concepts
\newcommand{\concept}[1]{\textbf{\color{darkblue}#1}}

\begin{document}

% Title slide
\begin{frame}
\titlepage
\end{frame}

% Reset section counter for this appendix
\setcounter{section}{0}

% Define sections for this appendix
\section{Introduction to HDLs}
\section{Chisel Language}
\section{SystemVerilog}
\section{Chisel vs SystemVerilog}
\section{Practical Examples}
\section{Design Methodology}

% Table of contents
\begin{frame}{Outline}
\tableofcontents
\end{frame}

% Learning objectives
\begin{frame}{Learning Objectives}
By the end of this appendix, you will understand:
\begin{itemize}
\item Hardware Description Language (HDL) fundamentals
\item Chisel language features and syntax
\item SystemVerilog capabilities and applications
\item Comparative analysis of Chisel vs SystemVerilog
\item Practical implementation examples in both languages
\item Design methodology and best practices
\item Integration with hardware development flows
\end{itemize}
\end{frame}

% Section 1: Introduction to HDLs
\section{Introduction to HDLs}

\begin{frame}{What are Hardware Description Languages?}
\begin{columns}
\begin{column}{0.6\textwidth}
\concept{HDL Definition}:
\begin{itemize}
\item Specialized programming languages for hardware design
\item Describe digital circuits at various abstraction levels
\item Enable simulation, synthesis, and verification
\item Bridge between algorithm and silicon implementation
\end{itemize}

\vspace{0.3cm}
\concept{Abstraction Levels}:
\begin{itemize}
\item \textbf{Behavioral}: Algorithm-level description
\item \textbf{RTL}: Register Transfer Level
\item \textbf{Gate}: Logic gate level
\item \textbf{Switch}: Transistor level
\end{itemize}

\vspace{0.3cm}
\concept{Key Benefits}:
\begin{itemize}
\item Design reusability and modularity
\item Simulation before fabrication
\item Automated synthesis and optimization
\item Verification and testing capabilities
\end{itemize}
\end{column}
\begin{column}{0.4\textwidth}
\begin{center}
\textbf{HDL Design Flow}
\vspace{0.3cm}

\tikz{
\node[draw, rectangle, minimum width=2cm, minimum height=0.8cm] (spec) at (0,4) {Specification};
\node[draw, rectangle, minimum width=2cm, minimum height=0.8cm] (hdl) at (0,3) {HDL Code};
\node[draw, rectangle, minimum width=2cm, minimum height=0.8cm] (sim) at (0,2) {Simulation};
\node[draw, rectangle, minimum width=2cm, minimum height=0.8cm] (syn) at (0,1) {Synthesis};
\node[draw, rectangle, minimum width=2cm, minimum height=0.8cm] (impl) at (0,0) {Implementation};

\draw[->] (spec) -- (hdl);
\draw[->] (hdl) -- (sim);
\draw[->] (hdl) -- (syn);
\draw[->] (syn) -- (impl);
}
\end{center}
\end{column}
\end{columns}
\end{frame}

\begin{frame}{Evolution of HDLs}
\begin{columns}
\begin{column}{0.6\textwidth}
\concept{Historical Development}:
\begin{itemize}
\item \textbf{1980s}: VHDL and Verilog emergence
\item \textbf{1990s}: Standardization and adoption
\item \textbf{2000s}: SystemVerilog and SystemC
\item \textbf{2010s}: High-level synthesis and Chisel
\item \textbf{2020s}: AI-assisted design and verification
\end{itemize}

\vspace{0.3cm}
\concept{Modern Trends}:
\begin{itemize}
\item Higher abstraction levels
\item Functional programming paradigms
\item Parameterizable and generic designs
\item Integration with software development tools
\item Open-source hardware ecosystems
\end{itemize}

\vspace{0.3cm}
\concept{Industry Adoption}:
\begin{itemize}
\item Traditional: Verilog/VHDL dominance
\item Emerging: Chisel, SpinalHDL, PyMTL
\item Specialized: HLS tools (Vivado HLS, Catapult)
\end{itemize}
\end{column}
\begin{column}{0.4\textwidth}
\begin{center}
\textbf{HDL Timeline}
\vspace{0.3cm}

\tikz{
\draw[thick] (0,0) -- (0,4);
\foreach \y/\year/\tech in {0/1980s/VHDL, 1/1990s/Verilog, 2/2000s/SystemV, 3/2010s/Chisel, 4/2020s/AI-HDL} {
  \draw[thick] (-0.1,\y) -- (0.1,\y);
  \node[left] at (-0.2,\y) {\year};
  \node[right] at (0.2,\y) {\tech};
}
}
\end{center>
\end{column>
\end{columns>
\end{frame>

% Section 2: Chisel Language
\section{Chisel Language}

\begin{frame}{Chisel Overview}
\begin{columns}
\begin{column}{0.6\textwidth}
\concept{Chisel Definition}:
\begin{itemize}
\item \textbf{C}onstructing \textbf{H}ardware \textbf{I}n a \textbf{S}cala \textbf{E}mbedded \textbf{L}anguage
\item Hardware design language embedded in Scala
\item Developed at UC Berkeley
\item Compiles to synthesizable Verilog
\end{itemize}

\vspace{0.3cm}
\concept{Key Features}:
\begin{itemize}
\item \textbf{Hardware Construction Language}: High-level abstractions
\item \textbf{Parameterization}: Generic and reusable modules
\item \textbf{Strong Typing}: Rigorous type checking
\item \textbf{Functional Programming}: Scala's functional paradigms
\item \textbf{Object-Oriented}: Classes and inheritance
\end{itemize>

\vspace{0.3cm}
\concept{Advantages}:
\begin{itemize}
\item Reduced design time and errors
\item Better code reusability
\item Integration with software tools
\item Modern programming language features
\end{itemize>
\end{column>
\begin{column}{0.4\textwidth>
\begin{center>
\textbf{Chisel Ecosystem}
\vspace{0.3cm}

\tikz{
\node[draw, circle, minimum width=1.5cm] (chisel) at (0,0) {Chisel};
\node[draw, rectangle] (scala) at (0,1.5) {Scala};
\node[draw, rectangle] (verilog) at (-1.5,-1) {Verilog};
\node[draw, rectangle] (firrtl) at (1.5,-1) {FIRRTL};
\node[draw, rectangle] (tools) at (0,-2) {CAD Tools};

\draw[->] (scala) -- (chisel);
\draw[->] (chisel) -- (verilog);
\draw[->] (chisel) -- (firrtl);
\draw[->] (verilog) -- (tools);
\draw[->] (firrtl) -- (tools);
}
\end{center>
\end{column>
\end{columns>
\end{frame>

\begin{frame}[fragile]{Chisel Basic Syntax}
\begin{columns}
\begin{column}{0.5\textwidth}
\concept{Import Statements}:
\begin{lstlisting}[style=chisel]
import chisel3._
import chisel3.util._
\end{lstlisting}

\vspace{0.3cm}
\concept{Basic Module}:
\begin{lstlisting}[style=chisel]
class MyModule extends Module {
  val io = IO(new Bundle {
    val in = Input(UInt(8.W))
    val out = Output(UInt(8.W))
  })
  
  io.out := io.in + 1.U
}
\end{lstlisting}
\end{column}
\begin{column}{0.5\textwidth}
\concept{Scala Basics for Chisel}:
\begin{itemize}
\item \texttt{val}: Immutable variables
\item \texttt{var}: Mutable variables  
\item \texttt{def}: Function definitions
\item \texttt{class}: Class definitions
\item \texttt{object}: Singleton objects
\end{itemize}

\vspace{0.3cm}
\concept{Chisel-Specific}:
\begin{itemize}
\item \texttt{Module}: Hardware module base class
\item \texttt{Bundle}: Interface definition
\item \texttt{IO()}: Port declaration
\item \texttt{:=}: Connection operator
\item \texttt{.W}: Width specification
\end{itemize>
\end{column>
\end{columns>
\end{frame>

\begin{frame}[fragile]{Chisel Data Types}
\begin{columns}
\begin{column}{0.5\textwidth}
\concept{Basic Types}:
\begin{lstlisting}[style=chisel]
// Unsigned integers
val a = UInt(8.W)
val b = 42.U

// Signed integers  
val c = SInt(16.W)
val d = (-5).S

// Boolean
val e = Bool()
val f = true.B
\end{lstlisting}

\vspace{0.3cm}
\concept{Aggregate Types}:
\begin{lstlisting}[style=chisel]
// Vector
val vec = Vec(4, UInt(8.W))

// Bundle
class MyBundle extends Bundle {
  val data = UInt(32.W)
  val valid = Bool()
}
\end{lstlisting>
\end{column>
\begin{column>{0.5\textwidth}
\concept{Type Features}:
\begin{itemize}
\item \textbf{Width inference}: Automatic width calculation
\item \textbf{Type safety}: Compile-time type checking
\item \textbf{Hardware semantics}: Types map to hardware
\end{itemize}

\vspace{0.3cm}
\concept{Width Specification}:
\begin{lstlisting}[style=chisel]
UInt(8.W)    // 8-bit unsigned
SInt(16.W)   // 16-bit signed
UInt()       // Inferred width
\end{lstlisting}

\vspace{0.3cm}
\concept{Literals}:
\begin{lstlisting}[style=chisel]
42.U         // Unsigned literal
(-5).S       // Signed literal
true.B       // Boolean literal
"h1A".U      // Hex literal
"b1010".U    // Binary literal
\end{lstlisting>
\end{column>
\end{columns>
\end{frame>

\begin{frame}[fragile]{Chisel Control Structures}
\begin{columns}
\begin{column}{0.5\textwidth}
\concept{Conditional Logic}:
\begin{lstlisting}[style=chisel]
when(condition) {
  // Execute when true
}.elsewhen(condition2) {
  // Execute when condition2 true
}.otherwise {
  // Default case
}
\end{lstlisting}

\vspace{0.3cm}
\concept{Switch Statement}:
\begin{lstlisting}[style=chisel]
switch(selector) {
  is(0.U) {
    // Case 0
  }
  is(1.U) {
    // Case 1
  }
}
\end{lstlisting>
\end{column}
\begin{column>{0.5\textwidth}
\concept{Registers}:
\begin{lstlisting}[style=chisel]
val reg = RegNext(io.in)
val regInit = RegInit(0.U(8.W))

// Clock and reset
val reg = withClock(clock) {
  withReset(reset) {
    Reg(UInt(8.W))
  }
}
\end{lstlisting}

\vspace{0.3cm}
\concept{Memory}:
\begin{lstlisting}[style=chisel]
val mem = Mem(1024, UInt(32.W))
val data = mem(addr)
mem(addr) := writeData
\end{lstlisting>
\end{column>
\end{columns>
\end{frame>

% Section 3: SystemVerilog
\section{SystemVerilog}

\begin{frame}{SystemVerilog Overview}
\begin{columns}
\begin{column}{0.6\textwidth}
\concept{SystemVerilog Definition}:
\begin{itemize}
\item Extension of Verilog-2001
\item IEEE 1800 standard
\item Combines design and verification features
\item Industry-standard HDL
\end{itemize}

\vspace{0.3cm}
\concept{Key Enhancements over Verilog}:
\begin{itemize}
\item \textbf{Object-Oriented Programming}: Classes and objects
\item \textbf{Advanced Data Types}: Structures, unions, enums
\item \textbf{Interfaces}: Modular connection specification
\item \textbf{Assertions}: Formal verification constructs
\item \textbf{Constrained Random}: Advanced testbench features
\end{itemize>

\vspace{0.3cm}
\concept{Application Areas}:
\begin{itemize}
\item RTL design and synthesis
\item Verification and testbenches
\item System-level modeling
\item Formal verification
\end{itemize>
\end{column>
\begin{column>{0.4\textwidth>
\begin{center>
\textbf{SystemVerilog Features}
\vspace{0.3cm}

\tikz{
\node[draw, circle, minimum width=2cm] (sv) at (0,0) {SystemVerilog};
\node[draw, rectangle] (design) at (-1.5,1.5) {Design};
\node[draw, rectangle] (verif) at (1.5,1.5) {Verification};
\node[draw, rectangle] (oop) at (-1.5,-1.5) {OOP};
\node[draw, rectangle] (assert) at (1.5,-1.5) {Assertions};

\draw[->] (sv) -- (design);
\draw[->] (sv) -- (verif);
\draw[->] (sv) -- (oop);
\draw[->] (sv) -- (assert);
}
\end{center>
\end{column>
\end{columns>
\end{frame>

\begin{frame}[fragile]{SystemVerilog Basic Syntax}
\begin{columns}
\begin{column}{0.5\textwidth}
\concept{Module Definition}:
\begin{lstlisting}[style=verilog]
module counter #(
  parameter WIDTH = 8
)(
  input  logic clk,
  input  logic rst_n,
  input  logic enable,
  output logic [WIDTH-1:0] count
);

always_ff @(posedge clk or negedge rst_n) begin
  if (!rst_n)
    count <= '0;
  else if (enable)
    count <= count + 1;
end

endmodule
\end{lstlisting>
\end{column}
\begin{column>{0.5\textwidth}
\concept{Key Syntax Elements}:
\begin{itemize}
\item \texttt{logic}: 4-state data type
\item \texttt{always\_ff}: Clocked logic
\item \texttt{always\_comb}: Combinational logic
\item \texttt{interface}: Connection specification
\item \texttt{modport}: Interface direction
\end{itemize>

\vspace{0.3cm}
\concept{Data Types}:
\begin{lstlisting}[style=verilog]
logic [7:0] data;
bit [15:0] address;
int counter;
real frequency;

typedef struct {
  logic [31:0] data;
  logic valid;
} packet_t;
\end{lstlisting>
\end{column>
\end{columns>
\end{frame>

\begin{frame}[fragile]{SystemVerilog Advanced Features}
\begin{columns}
\begin{column}{0.5\textwidth}
\concept{Interfaces}:
\begin{lstlisting}[style=verilog]
interface bus_if;
  logic [31:0] data;
  logic valid;
  logic ready;
  
  modport master (
    output data, valid,
    input ready
  );
  
  modport slave (
    input data, valid,
    output ready
  );
endinterface
\end{lstlisting}
\end{column}
\begin{column>{0.5\textwidth}
\concept{Classes (for Verification)}:
\begin{lstlisting}[style=verilog]
class packet;
  rand bit [31:0] data;
  rand bit [7:0] id;
  
  constraint valid_id {
    id inside {[1:10]};
  }
  
  function void display();
    $display("Packet: data=%h, id=%d", 
             data, id);
  endfunction
endclass
\end{lstlisting>

\vspace{0.3cm}
\concept{Assertions}:
\begin{lstlisting}[style=verilog]
assert property (
  @(posedge clk) 
  req |-> ##[1:3] ack
);
\end{lstlisting>
\end{column>
\end{columns>
\end{frame>

% Section 4: Chisel vs SystemVerilog
\section{Chisel vs SystemVerilog}

\begin{frame}{Comparative Analysis: Design Philosophy}
\begin{columns}
\begin{column}{0.5\textwidth}
\concept{Chisel Approach}:
\begin{itemize}
\item \textbf{Embedded DSL}: Leverages Scala ecosystem
\item \textbf{Functional Programming}: Immutable data, higher-order functions
\item \textbf{Generator-based}: Parameterizable hardware generators
\item \textbf{Type Safety}: Strong static typing
\item \textbf{Modern Tooling}: IDE support, debugging
\end{itemize>

\vspace{0.3cm}
\concept{Advantages}:
\begin{itemize}
\item Reduced boilerplate code
\item Better abstraction and reusability
\item Integration with software development
\item Powerful parameterization
\item Less error-prone
\end{itemize>
\end{column}
\begin{column>{0.5\textwidth}
\concept{SystemVerilog Approach}:
\begin{itemize}
\item \textbf{Traditional HDL}: Hardware-centric syntax
\item \textbf{Imperative Programming}: Procedural constructs
\item \textbf{Module-based}: Hierarchical design
\item \textbf{Industry Standard}: Wide tool support
\item \textbf{Verification Focus}: Advanced testbench features
\end{itemize>

\vspace{0.3cm}
\concept{Advantages}:
\begin{itemize}
\item Industry-wide adoption
\item Mature tool ecosystem
\item Direct hardware mapping
\item Comprehensive verification features
\item Familiar to hardware engineers
\end{itemize>
\end{column>
\end{columns>
\end{frame>

\begin{frame}[fragile]{Code Comparison: Simple Counter}
\begin{columns}
\begin{column}{0.5\textwidth}
\concept{Chisel Implementation}:
\begin{lstlisting}[style=chisel, basicstyle=\ttfamily\tiny]
class Counter(width: Int) extends Module {
  val io = IO(new Bundle {
    val enable = Input(Bool())
    val count = Output(UInt(width.W))
  })
  
  val reg = RegInit(0.U(width.W))
  
  when(io.enable) {
    reg := reg + 1.U
  }
  
  io.count := reg
}

// Usage
val counter8 = Module(new Counter(8))
val counter16 = Module(new Counter(16))
\end{lstlisting}
\end{column}
\begin{column>{0.5\textwidth}
\concept{SystemVerilog Implementation}:
\begin{lstlisting}[style=verilog, basicstyle=\ttfamily\tiny]
module counter #(
  parameter WIDTH = 8
)(
  input  logic clk,
  input  logic rst_n,
  input  logic enable,
  output logic [WIDTH-1:0] count
);

always_ff @(posedge clk or negedge rst_n) begin
  if (!rst_n)
    count <= '0;
  else if (enable)
    count <= count + 1;
end

endmodule

// Usage
counter #(.WIDTH(8)) counter8_inst (...);
counter #(.WIDTH(16)) counter16_inst (...);
\end{lstlisting>
\end{column>
\end{columns>
\end{frame>

\begin{frame}{Performance and Productivity Comparison}
\begin{columns}
\begin{column}{0.6\textwidth}
\concept{Development Productivity}:
\begin{itemize}
\item \textbf{Chisel}: Higher abstraction reduces code size by 2-10x
\item \textbf{SystemVerilog}: More verbose but explicit
\item \textbf{Learning Curve}: Chisel requires Scala knowledge
\item \textbf{Debugging}: Chisel provides better error messages
\end{itemize>

\vspace{0.3cm}
\concept{Generated Hardware Quality}:
\begin{itemize}
\item \textbf{Performance}: Comparable synthesized results
\item \textbf{Area}: Similar resource utilization
\item \textbf{Power}: No significant difference
\item \textbf{Timing}: Both meet timing constraints equally
\end{itemize>

\vspace{0.3cm}
\concept{Tool Ecosystem}:
\begin{itemize}
\item \textbf{SystemVerilog}: Universal tool support
\item \textbf{Chisel}: Growing but limited tool support
\item \textbf{Verification}: SystemVerilog has mature frameworks
\item \textbf{Simulation}: Both integrate with standard simulators
\end{itemize>
\end{column>
\begin{column>{0.4\textwidth>
\begin{center>
\textbf{Comparison Metrics}
\vspace{0.3cm}

\tikz{
% Code size comparison
\draw[thick, fill=blue!30] (0,3.5) rectangle (1,4);
\draw[thick, fill=red!30] (0,3) rectangle (2.5,3.5);
\node at (-0.5,3.75) {Chisel};
\node at (-0.5,3.25) {SystemV};
\node at (1.5,4.2) {\textbf{Code Size}};

% Learning curve
\draw[thick, fill=red!30] (0,2.5) rectangle (2,3);
\draw[thick, fill=blue!30] (0,2) rectangle (1.2,2.5);
\node at (-0.5,2.75) {Chisel};
\node at (-0.5,2.25) {SystemV};
\node at (1.5,3.2) {\textbf{Learning}};

% Tool support
\draw[thick, fill=blue!30] (0,1.5) rectangle (1.5,2);
\draw[thick, fill=red!30] (0,1) rectangle (2.8,1.5);
\node at (-0.5,1.75) {Chisel};
\node at (-0.5,1.25) {SystemV};
\node at (1.5,2.2) {\textbf{Tools}};
}
\end{center>
\end{column>
\end{columns>
\end{frame>

% Section 5: Practical Examples
\section{Practical Examples}

\begin{frame}[fragile]{Example: FIFO Implementation in Chisel}
\begin{columns}
\begin{column}{0.5\textwidth}
\begin{lstlisting}[style=chisel, basicstyle=\ttfamily\tiny]
class FIFO(depth: Int, width: Int) extends Module {
  val io = IO(new Bundle {
    val enq = Flipped(new DecoupledIO(UInt(width.W)))
    val deq = new DecoupledIO(UInt(width.W))
  })
  
  val mem = Mem(depth, UInt(width.W))
  val enqPtr = RegInit(0.U(log2Ceil(depth).W))
  val deqPtr = RegInit(0.U(log2Ceil(depth).W))
  val full = RegInit(false.B)
  val empty = RegInit(true.B)
  
  // Enqueue logic
  when(io.enq.fire) {
    mem(enqPtr) := io.enq.bits
    enqPtr := enqPtr + 1.U
  }
  
  // Dequeue logic
  when(io.deq.fire) {
    deqPtr := deqPtr + 1.U
  }
  
  io.enq.ready := !full
  io.deq.valid := !empty
  io.deq.bits := mem(deqPtr)
}
\end{lstlisting>
\end{column>
\begin{column>{0.5\textwidth}
\concept{Key Chisel Features Demonstrated}:
\begin{itemize}
\item \textbf{Parameterization}: Generic depth and width
\item \textbf{Bundle}: Interface definition
\item \textbf{DecoupledIO}: Ready-valid interface
\item \textbf{Mem}: Memory primitive
\item \textbf{RegInit}: Initialized registers
\item \textbf{fire}: Ready and valid both true
\end{itemize>

\vspace{0.3cm}
\concept{Design Benefits}:
\begin{itemize}
\item Type-safe interfaces
\item Automatic width inference
\item Reusable across projects
\item Clear separation of concerns
\item Built-in flow control
\end{itemize>
\end{column>
\end{columns>
\end{frame>

\begin{frame}[fragile]{Example: FIFO Implementation in SystemVerilog}
\begin{columns}
\begin{column}{0.5\textwidth}
\begin{lstlisting}[style=verilog, basicstyle=\ttfamily\tiny]
module fifo #(
  parameter DEPTH = 16,
  parameter WIDTH = 8,
  parameter ADDR_WIDTH = $clog2(DEPTH)
)(
  input  logic clk,
  input  logic rst_n,
  input  logic [WIDTH-1:0] din,
  input  logic wr_en,
  input  logic rd_en,
  output logic [WIDTH-1:0] dout,
  output logic full,
  output logic empty
);

logic [WIDTH-1:0] mem [DEPTH-1:0];
logic [ADDR_WIDTH-1:0] wr_ptr, rd_ptr;
logic [ADDR_WIDTH:0] count;

always_ff @(posedge clk or negedge rst_n) begin
  if (!rst_n) begin
    wr_ptr <= '0;
    rd_ptr <= '0;
    count <= '0;
  end else begin
    if (wr_en && !full) begin
      mem[wr_ptr] <= din;
      wr_ptr <= wr_ptr + 1;
    end
    if (rd_en && !empty) begin
      rd_ptr <= rd_ptr + 1;
    end
    
    case ({wr_en && !full, rd_en && !empty})
      2'b10: count <= count + 1;
      2'b01: count <= count - 1;
    endcase
  end
end

assign dout = mem[rd_ptr];
assign full = (count == DEPTH);
assign empty = (count == 0);

endmodule
\end{lstlisting>
\end{column>
\begin{column>{0.5\textwidth}
\concept{SystemVerilog Features Demonstrated}:
\begin{itemize}
\item \textbf{Parameters}: Configurable design
\item \textbf{Memory Array}: Explicit memory declaration
\item \textbf{always\_ff}: Clocked logic
\item \textbf{case}: Multi-way selection
\item \textbf{System Functions}: \$clog2 for width calculation
\end{itemize>

\vspace{0.3cm}
\concept{Implementation Details}:
\begin{itemize}
\item Explicit pointer management
\item Manual count tracking
\item Direct memory indexing
\item Combinational output assignment
\item Reset handling
\end{itemize>
\end{column>
\end{columns>
\end{frame>

\begin{frame}[fragile]{Example: Parameterizable ALU}
\begin{columns}
\begin{column}{0.5\textwidth}
\concept{Chisel ALU}:
\begin{lstlisting}[style=chisel, basicstyle=\ttfamily\tiny]
object ALUOp extends ChiselEnum {
  val ADD, SUB, AND, OR, XOR = Value
}

class ALU(width: Int) extends Module {
  val io = IO(new Bundle {
    val a = Input(UInt(width.W))
    val b = Input(UInt(width.W))
    val op = Input(ALUOp())
    val result = Output(UInt(width.W))
    val zero = Output(Bool())
  })
  
  io.result := MuxLookup(io.op, 0.U)(Seq(
    ALUOp.ADD -> (io.a + io.b),
    ALUOp.SUB -> (io.a - io.b),
    ALUOp.AND -> (io.a & io.b),
    ALUOp.OR  -> (io.a | io.b),
    ALUOp.XOR -> (io.a ^ io.b)
  ))
  
  io.zero := io.result === 0.U
}
\end{lstlisting>
\end{column>
\begin{column>{0.5\textwidth}
\concept{SystemVerilog ALU}:
\begin{lstlisting}[style=verilog, basicstyle=\ttfamily\tiny]
typedef enum logic [2:0] {
  ADD = 3'b000,
  SUB = 3'b001,
  AND = 3'b010,
  OR  = 3'b011,
  XOR = 3'b100
} alu_op_t;

module alu #(
  parameter WIDTH = 32
)(
  input  logic [WIDTH-1:0] a,
  input  logic [WIDTH-1:0] b,
  input  alu_op_t op,
  output logic [WIDTH-1:0] result,
  output logic zero
);

always_comb begin
  case (op)
    ADD: result = a + b;
    SUB: result = a - b;
    AND: result = a & b;
    OR:  result = a | b;
    XOR: result = a ^ b;
    default: result = '0;
  endcase
end

assign zero = (result == '0);

endmodule
\end{lstlisting>
\end{column>
\end{columns>
\end{frame>

% Section 6: Design Methodology
\section{Design Methodology}

\begin{frame}{HDL Design Best Practices}
\begin{columns}
\begin{column}{0.6\textwidth}
\concept{General Design Principles}:
\begin{itemize}
\item \textbf{Modularity}: Break design into reusable components
\item \textbf{Hierarchy}: Use clear hierarchical structure
\item \textbf{Parameterization}: Make designs configurable
\item \textbf{Documentation}: Comment code and interfaces
\item \textbf{Testing}: Comprehensive verification strategy
\end{itemize}

\vspace{0.3cm}
\concept{Chisel-Specific Best Practices}:
\begin{itemize}
\item Use Bundle for complex interfaces
\item Leverage Scala's type system
\item Create hardware generators, not just modules
\item Use functional programming constructs
\item Separate generation-time and runtime logic
\end{itemize>

\vspace{0.3cm}
\concept{SystemVerilog Best Practices}:
\begin{itemize}
\item Use interfaces for complex connections
\item Employ proper coding styles for synthesis
\item Separate design and verification code
\item Use assertions for specification
\item Follow naming conventions
\end{itemize>
\end{column>
\begin{column>{0.4\textwidth>
\begin{center>
\textbf{Design Flow}
\vspace{0.3cm}

\tikz{
\node[draw, rectangle, minimum width=2cm, minimum height=0.6cm] (spec) at (0,4) {Specification};
\node[draw, rectangle, minimum width=2cm, minimum height=0.6cm] (arch) at (0,3.3) {Architecture};
\node[draw, rectangle, minimum width=2cm, minimum height=0.6cm] (hdl) at (0,2.6) {HDL Coding};
\node[draw, rectangle, minimum width=2cm, minimum height=0.6cm] (test) at (0,1.9) {Testing};
\node[draw, rectangle, minimum width=2cm, minimum height=0.6cm] (syn) at (0,1.2) {Synthesis};
\node[draw, rectangle, minimum width=2cm, minimum height=0.6cm] (impl) at (0,0.5) {Implementation};

\draw[->] (spec) -- (arch);
\draw[->] (arch) -- (hdl);
\draw[->] (hdl) -- (test);
\draw[->] (test) -- (syn);
\draw[->] (syn) -- (impl);

% Feedback loops
\draw[->, dashed] (test.west) -- ++(-0.5,0) -- ++(0,0.7) -- (hdl.west);
\draw[->, dashed] (syn.west) -- ++(-0.7,0) -- ++(0,1.4) -- (hdl.west);
}
\end{center>
\end{column>
\end{columns>
\end{frame>

\begin{frame}{Tool Integration and Workflow}
\begin{columns}
\begin{column}{0.6\textwidth}
\concept{Chisel Workflow}:
\begin{itemize}
\item \textbf{Development}: IntelliJ IDEA or VS Code with Scala
\item \textbf{Build System}: SBT (Scala Build Tool)
\item \textbf{Testing}: ChiselTest framework
\item \textbf{Generation}: FIRRTL compiler to Verilog
\item \textbf{Synthesis}: Standard EDA tools (Synopsys, Cadence)
\end{itemize>

\vspace{0.3cm}
\concept{SystemVerilog Workflow}:
\begin{itemize}
\item \textbf{Development}: Specialized HDL editors
\item \textbf{Simulation}: ModelSim, VCS, Xcelium
\item \textbf{Synthesis}: Design Compiler, Genus
\item \textbf{Verification}: UVM, formal tools
\item \textbf{Implementation}: Place and route tools
\end{itemize>

\vspace{0.3cm}
\concept{Integration Considerations}:
\begin{itemize}
\item Version control strategies
\item Continuous integration setup
\item Documentation generation
\item IP reuse and packaging
\end{itemize>
\end{column>
\begin{column>{0.4\textwidth>
\begin{center>
\textbf{Tool Ecosystem}
\vspace{0.3cm}

\tikz{
% Chisel tools
\node at (0,3.5) {\textbf{Chisel}};
\draw[thick] (-1,3) rectangle (1,3.3);
\node at (0,3.15) {SBT/IDE};
\draw[thick] (-1,2.6) rectangle (1,2.9);
\node at (0,2.75) {ChiselTest};
\draw[thick] (-1,2.2) rectangle (1,2.5);
\node at (0,2.35) {FIRRTL};

% SystemVerilog tools
\node at (0,1.8) {\textbf{SystemVerilog}};
\draw[thick] (-1,1.3) rectangle (1,1.6);
\node at (0,1.45) {Simulators};
\draw[thick] (-1,0.9) rectangle (1,1.2);
\node at (0,1.05) {Synthesis};
\draw[thick] (-1,0.5) rectangle (1,0.8);
\node at (0,0.65) {Verification};
}
\end{center>
\end{column>
\end{columns>
\end{frame>

% Summary
\section{Summary}

\begin{frame}{Appendix Summary}
\begin{columns}
\begin{column}{0.6\textwidth}
\concept{Key Takeaways}:
\begin{itemize}
\item HDLs are essential for modern hardware design
\item Chisel offers higher abstraction and productivity
\item SystemVerilog provides industry-standard features
\item Both languages produce equivalent hardware
\item Choice depends on project requirements and team expertise
\item Tool ecosystem maturity varies between languages
\end{itemize>

\vspace{0.3cm}
\concept{When to Choose Chisel}:
\begin{itemize}
\item Complex, parameterizable designs
\item Research and academic projects
\item Teams with software background
\item Need for rapid prototyping
\end{itemize>

\vspace{0.3cm}
\concept{When to Choose SystemVerilog}:
\begin{itemize}
\item Industry production designs
\item Existing SystemVerilog codebase
\item Comprehensive verification requirements
\item Maximum tool compatibility needed
\end{itemize>
\end{column>
\begin{column>{0.4\textwidth>
\begin{center>
\textbf{HDL Selection Criteria}
\vspace{0.3cm}

\tikz{
% Decision factors
\node at (1,3.5) {\textbf{Project Factors}};

\draw[thick] (0,3) rectangle (2,3.2);
\node at (1,3.1) {Team Expertise};

\draw[thick] (0,2.6) rectangle (2,2.8);
\node at (1,2.7) {Tool Requirements};

\draw[thick] (0,2.2) rectangle (2,2.4);
\node at (1,2.3) {Design Complexity};

\draw[thick] (0,1.8) rectangle (2,2);
\node at (1,1.9) {Time Constraints};

\draw[thick] (0,1.4) rectangle (2,1.6);
\node at (1,1.5) {Verification Needs};

% Decision arrow
\draw[->, thick] (1,1.2) -- (1,0.8);
\node at (1,0.6) {\textbf{HDL Choice}};
}
\end{center>
\end{column>
\end{columns>
\end{frame>

\begin{frame}{Further Resources}
\begin{columns}
\begin{column}{0.6\textwidth}
\concept{Chisel Resources}:
\begin{itemize}
\item \textbf{Official Documentation}: chisel-lang.org
\item \textbf{Book}: "Digital Design with Chisel" (free online)
\item \textbf{Cheat Sheet}: Quick reference guide
\item \textbf{GitHub}: chisel-lang/chisel3
\item \textbf{Community}: Gitter chat, Stack Overflow
\end{itemize>

\vspace{0.3cm}
\concept{SystemVerilog Resources}:
\begin{itemize}
\item \textbf{IEEE Standard}: 1800-2017
\item \textbf{Books}: "SystemVerilog for Design" by Sutherland
\item \textbf{Training}: Verification Academy
\item \textbf{Forums}: EDA community forums
\item \textbf{Tools}: Vendor documentation
\end{itemize>

\vspace{0.3cm}
\concept{Practice Recommendations}:
\begin{itemize}
\item Start with simple examples
\item Build incrementally complex designs
\item Focus on one language initially
\item Practice with real hardware targets
\end{itemize>
\end{column>
\begin{column>{0.4\textwidth>
\begin{center>
\textbf{Learning Path}
\vspace{0.3cm}

\tikz{
\node[draw, rectangle, minimum width=2cm, minimum height=0.6cm] (basics) at (0,3) {HDL Basics};
\node[draw, rectangle, minimum width=2cm, minimum height=0.6cm] (syntax) at (0,2.3) {Language Syntax};
\node[draw, rectangle, minimum width=2cm, minimum height=0.6cm] (examples) at (0,1.6) {Simple Examples};
\node[draw, rectangle, minimum width=2cm, minimum height=0.6cm] (projects) at (0,0.9) {Real Projects};
\node[draw, rectangle, minimum width=2cm, minimum height=0.6cm] (advanced) at (0,0.2) {Advanced Topics};

\draw[->] (basics) -- (syntax);
\draw[->] (syntax) -- (examples);
\draw[->] (examples) -- (projects);
\draw[->] (projects) -- (advanced);
}
\end{center>
\end{column>
\end{columns>
\end{frame>

\end{document}

